% Physical unit 009; accepted Session 02.
\GordanUnitStart{009}{printed page 4}
\GordanTopFolio{4}
\begin{center}
{\large CAPUT II.}\\[0.8em]
\S\ 1.
\end{center}
Curva puncto $\Pi$ in sphaeroide descripta, nulla nisi normali vi movente, nominatur linea geodetica. Aequationem invenimus ejiciendo ex motus corporei puncti aequationibus quantitates $t$ et $\lambda$, quarum haec vim normalem, illa tempus proportionale habet.

\begin{center}\S\ 2.\end{center}
Aequatione igitur sphaeroidis \gq{$f=0$} sumta, aequationes motus in coordinatis rectangulis hae sunt:
\[
\begin{gathered}
 {d_2x\over dt^2}=\lambda{\partial f\over\partial x},\\[0.25em]
 {d_2y\over dt^2}=\lambda{\partial f\over\partial y},\\[0.25em]
 {d_2z\over dt^2}=\lambda{\partial f\over\partial z}.
\end{gathered}
\]
quae si deinceps differentialibus
\[
 dx\qquad dy\qquad dz
\]
multiplicatae adduntur; fit:
\[
 {dx\over dt}{d_2x\over dt^2}+{dy\over dt}{d_2y\over dt^2}+{dz\over dt}{d_2z\over dt^2}=0
\]
vel integrando:
\[
 {dx\over dt}{dx\over dt}+{dy\over dt}{dy\over dt}+{dz\over dt}{dz\over dt}={1\over C^2}
\]
sive:
\[
 dt=C\,ds,
\]
unde videmus notum illud, celeritatem puncti $\Pi$ esse constantem.

\begin{center}\S\ 3.\end{center}
Hoc valore in aequationes motus inducto, habemus:
\[
\begin{gathered}
 {d_2x\over ds^2}=\lambda C^2{\partial f\over\partial x},\\[0.25em]
 {d_2y\over ds^2}=\lambda C^2{\partial f\over\partial y},\\[0.25em]
 {d_2z\over ds^2}=\lambda C^2{\partial f\over\partial z}.
\end{gathered}
\]
quibus nihil aliud efficitur, nisi, ut variatio extensionis, qua in constituenda ratio $f=0$ teneatur, ipsa evanescat.

Qua re aequatio lineae geodeticae est:
